{# This is a Jinja2 template for a latex file.  Both use braces heavily! 
    We use the raw ... endraw statements to surround untouched latex. 
    We use the `embrace` filter to enclose an argument with braces. 
    Braces can also be quoted with the `bgroup` and `egroup` variables
    or with `{{ "{" }}` and  `{{ "}" }}`.
#}
{%- set bgroup = "{" -%}
{%- set egroup = "}" -%}
{%- macro random_int() %}{% for n in [0,1,2,3,4,5,6] %}{{ [0,1,2,3,4,5,6,7,8,9]|random }}{% endfor %}{% endmacro -%}
{%- if cookiecutter.has_versions -%}
{%- set versions = cookiecutter.versions_csv.split(',') | map('trim') | list -%}
{#-
  Determine which versions should have separate solution files.
  If the cookiecutter context provides `versions_with_solutions`, use that
  CSV; otherwise default to all versions (when versions is set).
-#}
{%- if cookiecutter.get('versions_with_solutions') -%}
{%- set versions_with_solutions = cookiecutter['versions_with_solutions'].split(',') -%}
{%- elif versions -%}
{%- set versions_with_solutions = versions -%}
{%- else -%}
{%- set versions_with_solutions = None -%}
{%- endif -%}
{%- else -%}
{%- set versions = None -%}
{%- endif -%}
%<*driver>
%%!TEX TS-program = dtxmake
%%!TEX dtxmake-subengine = lualatexmk
\input docstrip.tex
\askforoverwritefalse
\generate{
{%- if versions -%}
{%- for v in versions %}
  \file{\jobname-{{v}}.qns.tex}{\from{\jobname.dtx}{questions,{{v}}}}
  {%- if versions_with_solutions is none or (v in versions_with_solutions) %}
  \file{\jobname-{{v}}.sol.tex}{\from{\jobname.dtx}{questions,solutions,{{v}}}}
  {%- endif -%}
{%- endfor -%}
{%- else -%}
    \file{\jobname.qns.tex}{\from{\jobname.dtx}{questions}}
    \file{\jobname.sol.tex}{\from{\jobname.dtx}{questions,solutions}}
{%- endif %}
}
\endbatchfile
%</driver>
%<*questions>
\documentclass[addpoints
%<solutions>,answers
{%- if cookiecutter.use_nyu_fonts %}
,nyufonts
{%- endif %}
{%- if cookiecutter.get('use_barcode') %}
,barcode
{%- endif %}
{%- if cookiecutter.use_bubble_sheet %}
,nocirclechoices
{%- endif %}
]{{ bgroup }}nyuexam3{{ egroup }}
\ProvidesFile
%<*dtx>
{{ (cookiecutter.exam_code + ".dtx") | embrace }}
%</dtx>
{%- if versions -%}
{%- for v in versions %}
%<{{v}}&questions&!solutions>{{ (cookiecutter.exam_code + "-" + v + ".qns.tex") | embrace }}
{%- if versions_with_solutions is none or (v in versions_with_solutions) %}
%<{{v}}&questions&solutions>{{ (cookiecutter.exam_code + "-" + v + ".sol.tex") | embrace }}
{%- endif -%}
{%- endfor -%}
{%- else %}
%<questions&!solutions>{{ (cookiecutter.exam_code + ".qns.tex") | embrace }}
%<questions&solutions>{{ (cookiecutter.exam_code + ".sol.tex") | embrace }}
{%- endif %}
    [{% now 'local', '%Y/%m/%d' %} v0.0 {{cookiecutter.course_name}}, {{cookiecutter.term_name}}, {{cookiecutter.exam_name}}]
\title{{cookiecutter.exam_name | embrace }}
\author{{cookiecutter.instructor_name | replace('. ', '.~') | embrace }}
{%- if cookiecutter.get('version_times') %}
\date{{ bgroup }}{{ cookiecutter.exam_date | localize_date }}\relax
{%- for pair in cookiecutter.version_times.split(',') %}
{%- set v, t = pair.split(':', 1) %}
%<{{ v | trim }}>, {{ t | trim }}
{%- endfor %}
{{ egroup }}
{%- else %}
\date{{ cookiecutter.exam_date | localize_date | embrace }}
{%- endif %}
{%- set short_course_opt = '[' ~ cookiecutter.short_course_name ~ ']' if cookiecutter.get('short_course_name') and cookiecutter.short_course_name != cookiecutter.course_name else '' %}
\course{{ short_course_opt }}{{ cookiecutter.course_name | embrace }}
{%- if cookiecutter.get('version_course_names') -%}
{%- for group in cookiecutter.version_course_names.split(',') %}
{%- set gv, gname = group.split(':', 1) %}
%<{%- for v in gv.split(';') -%}{{ v | trim }}{% if not loop.last %}|{% endif %}{%- endfor -%}>\course{{ short_course_opt }}{{ gname | trim | embrace }}
{%- endfor -%}
{%- endif %}
{% raw %}

\usepackage{markversion}% Sets \VCRevisionMod and provides \UseVersionOf
\footer% all non-cover pages
    {\UseVersionOf{\jobname.tex}, \UseDateOf{\jobname.tex}}
    {\thepage/\numpages}
    {Revision \VCRevisionMod}
\makeatletter
\coverfooter% footer on cover pages only
    {\run@lfoot}
    {}
    {\run@rfoot}
\makeatother



\usepackage[overload]{exam-randomizechoices}
% Update the seed below to keep from copy & pasting the same one
% from exam to exam.
{% endraw -%}
% Random seed(s) generated {% now 'local', '%c' %} 
{%- if versions and cookiecutter.get('version_randomization_groups') -%}
{%- for g in cookiecutter.version_randomization_groups.split(',') %}
%<{%- for v in g.split(';') -%}{{v}}{% if not loop.last %}|{% endif %}{%- endfor -%}>\def\randomseed{{ random_int() | embrace }}
{%- endfor %}
% Shared by all versions, so \GetVersionItem gives each version a different item
\def\sharedseed{{ random_int() | embrace }}
{%- for g in cookiecutter.version_randomization_groups.split(',') %}
%<{%- for v in g.split(';') -%}{{v}}{% if not loop.last %}|{% endif %}{%- endfor -%}>\def\versionindex{{ loop.index | string | embrace }}
{%- endfor -%}
{%- else %}
\def\randomseed{{ random_int() | embrace }}
{%- endif %}
{% raw -%}
\setrandomizerseed{\randomseed}
\pgfmathsetseed{\randomseed}
{% endraw -%}
{#- The check-all-that-apply groups and true/false statements are shuffled with randomlist. -#}
{%- if cookiecutter.get('use_randomlist') or (cookiecutter.num_cata | int) > 0 or (cookiecutter.num_tf | int) > 0 -%}
\usepackage{randomlist}
\RLsetrandomseed{\randomseed}
\usepackage{versionlist}
{%- endif -%}
{%- raw %}
\ifluatex
  \directlua{math.randomseed(token.get_macro("randomseed"))}
  \usepackage{luacode}
  % See the luacode manual to understand why we need the \luaexec command here.
  % Should be single backlash, single percent, in double quotes
  \luaexec{PERCENT_CHAR = "\%"}
  \directlua{utils = require("utils")}
  % Write the answer key to \jobname.csv (question, answer, points)
  \ifprintanswers
    \directlua{aux_to_csv = require("aux_to_csv")}
    \directlua{aux_to_csv.register_for_job(tex.jobname)}
  \fi
\fi

{% endraw -%}
{%- if cookiecutter.get('extra_latex_packages') %}
{% set extra_packages = cookiecutter['extra_latex_packages'].split(',') %}
{% for pkg in extra_packages -%}
\usepackage{{ pkg | trim | embrace }}
{% endfor %}
{%- endif -%}
{%- if cookiecutter.get('extra_tikz_libraries') -%}
{% set extra_libraries = cookiecutter['extra_tikz_libraries'].split(',') -%}
{%- for lib in extra_libraries -%}
\usetikzlibrary{{ lib | trim | embrace }}
{% endfor -%}
{%- endif %}
{%- raw %}


\newcommand{\tf}[1][{}]{%
  \fillin[#1][2em]%
}

\setlength{\answerskip}{1cm}

% points of question types
\newcommand{\mcpoints}{
{%- endraw -%}
{{ cookiecutter.points_per_mcq }}
{%- raw -%}
}% multiple choice
\newcommand{\mspoints}{1}% each option of a check-all-that-apply group
\newcommand{\tfpoints}{1}% truth value of a true/false statement
\newcommand{\tfrpoints}{2}% reason for a true/false statement

\usepackage{amssymb}% \Box, on the cover
\usepackage{siunitx}% \qty command
{%- endraw %}
{%- if cookiecutter.get('use_formula_sheet') %}
\usepackage{pdfpages}% \includepdf, for the formula sheet
{%- endif -%}
{%- if cookiecutter.icons_package == 'marvosym' -%}
{%- raw %}
\let\Cross\relax
\usepackage{marvosym} %\Stopsign
\newcommand{\iconstop}{\Stopsign}
{% endraw -%}
{%- elif cookiecutter.icons_package == 'fontawesome5' -%}
{%- raw %}
\usepackage{fontawesome5}
\newcommand{\iconstop}{\faHandPaper[regular]}
{% endraw -%}
{%- else %}
% No icons package detected. Probably gonna get an error later on.
{% endif -%}
{%- raw -%}
\newcommand{\mystop}{\raisebox{-0.5\totalheight}{\fontsize{50}{60}\color{Magenta}\selectfont\iconstop}}

\input{command-verbs.tex}% glossary of command terms (explain, prove, ...)

\begin{document}

\begin{coverpages}
    \vspace*{\fill}
    \maketitle
    {\setlength{\parindent}{0pt}%
    \begin{minipage}[b]{8cm}
        \hfill\fbox{
            \begin{minipage}[b][1.35cm][t]{.9\linewidth}
              Name:
    
              \ifprintanswers\sffamily\Huge\color{UltraViolet} Solutions\fi
            \end{minipage}
      }\hfill\vspace{5ex}
    \end{minipage}
    \begin{minipage}[b]{8cm}
        \hfill\fbox{
            \begin{minipage}[b][1.35cm][t]{.9\linewidth}
              NetID:
            \end{minipage}
      }\hfill\vspace{5ex}
    \end{minipage}
    \hspace*{\fill}
    }
    
\bigskip
\begin{instructions}%
    \noindent{\textbf{\uppercase{READ THE FOLLOWING INFORMATION.}}}
    
    \begin{itemize}
        \item You have \textbf{
{%- endraw -%}
{{ cookiecutter.exam_duration }}{{ " " }} 
{%- raw -%}        
        minutes} to finish this exam.
        \item Fixed-response questions on this exam will be graded by computer.
        Please fill in the circles completely and use pencil in case you want to
        change your responses.  
        \item Calculators, books, notes, blank scrap paper, and other aids are not allowed.
        \item You may use the backs of the pages or the extra pages for scratch work.  
        \textbf{\emph{Do not unstaple or remove pages}}.
    \end{itemize}

    
    
    {\noindent\mystop
    \hfill \textbf{\uppercase{Do not begin this exam until signaled to do so.}}
    \hfill \mystop}
    
\end{instructions}

\vfill
\noindent\rule{\textwidth}{1pt}
\emph{For office use only}

\noindent\rule{0cm}{1cm}
E\# \rule{0.75in}{0.4pt}
S\# \rule{0.75in}{0.4pt}
CO: \hrulefill\ 
CI: \hrulefill\ 
Fin: \hrulefill\ 
FT \(\Box\)
\clearpage
\printglossaries
\label{command-verbs-glossary}
\clearpage
{% endraw -%}
{%- if cookiecutter.use_bubble_sheet -%}
{#- 
    Bubble rows used: one per MCQ, five per check-all group, two per T/F statement.
    The sheet has columns of 25 rows. Row spacing was calibrated on
    Fall 2026 Discrete Midterm I: 0.2375in per row, plus 0.25in between groups of five.
-#}
{%- set nrows = (cookiecutter.num_mcq | int) + 5 * (cookiecutter.num_cata | int) + 2 * (cookiecutter.num_tf | int) -%}
{%- if nrows <= 25 -%}
{%- set first_unused = nrows + 1 -%}
{%- else -%}
{%- set first_unused = nrows - 25 + 1 -%}
{%- endif -%}
{%- set ns = namespace(gaps=0) -%}
{%- for m in range(first_unused, 25) -%}{%- if m % 5 == 0 -%}{%- set ns.gaps = ns.gaps + 1 -%}{%- endif -%}{%- endfor -%}
{%- set whiteout_height = 0.06 + 0.2375 * (26 - first_unused) + 0.25 * ns.gaps -%}
{%- raw -%}
\ifprintanswers\else
\thispagestyle{empty}
\begin{tikzpicture}[
    overlay,
    remember picture,
    whiteout/.style={draw=none,fill=white},
    highlightbox/.style={draw=Magenta,line width=1.8pt,rounded corners=4pt}
]
    \node at (current page.center) {
{%  endraw -%}
{% set letters = ['A','B','C','D','E'] %}
{%- if versions -%}
{%- for v in versions -%}
%<{{v}}>\includegraphics{{ ("mc_bubble_sheet_" + letters[loop.index0 % 5]) | embrace}}
{% endfor -%}
{%- else %}
\includegraphics{{ ("mc_bubble_sheet_" + letters[0]) | embrace}}
{% endif %}
{%  raw -%}
};
    % White out unused rows to avoid confusion
    % LL = lower left corner of bubble matrix
    % UR = upper right corner of bubble matrix
    % the origin is the upper left corner of the page's printed area
    % (1 inch in each direction from upper left corner of the stock?)
    \coordinate (LL) at (-0.5in,-9in);
    \coordinate (UR) at (6.5in,-1.75in); 
    \coordinate (UL) at (LL |- UR);
    \coordinate (LR) at (LL -| UR);
    % highlight Name/NetID fields at the top of the bubble sheet
    \draw[highlightbox] ([xshift=0.12in,yshift=1.10in]UL) rectangle ++(3.4in,1.1in);
{%- endraw %}
    % {{ nrows }} bubble rows are used.
    % Heights are computed by the cookiecutter; check the rendered page and adjust if needed.
{%- if nrows <= 25 %}
    % white out columns 2-4
    \draw[whiteout] ([xshift=1.6in]LL) rectangle (UR);
    % white out rows {{ nrows + 1 }}--25 in column 1
    \draw[whiteout] (LL) rectangle ++(1.6in,{{ "%.2f" | format(whiteout_height) }}in);
{%- elif nrows < 50 %}
    % white out columns 3-4
    \draw[whiteout] ([xshift=3.4in]LL) rectangle (UR);
    % white out rows {{ nrows + 1 }}--50 in column 2
    \draw[whiteout] ([xshift=1.75in]LL) rectangle ++(1.75in,{{ "%.2f" | format(whiteout_height) }}in);
{%- elif nrows == 50 %}
    % white out columns 3-4
    \draw[whiteout] ([xshift=3.4in]LL) rectangle (UR);
{%- else %}
    % More than 50 rows: nothing whited out; consider adjusting by hand.
{%- endif %}
{%- raw %}
    \node[anchor=north west]
        at (3.25in,-3in)
        {\begin{minipage}{0.5\textwidth}
            \noindent{\textbf{\uppercase{Bubble Sheet Instructions}}}
            \begin{instructions}
            \begin{itemize}
                \item \textbf{\color{Magenta}Write your Name and NetID at the top of this bubble sheet.}
                \item Fill in your answers completely. Pencil is recommended.
                \item Erase your mistakes completely.
                \item The number of questions shown is the number that you should answer.
            \end{itemize}
            \end{instructions}
        \end{minipage}};
    \tikzset{
        bubble sheet fill in/.style={
            draw=none,
            font={\bfseries\sffamily\Large},
            DeepViolet
        }
    }
    \node[bubble sheet fill in,anchor=south west]
        at (3.6in,-0.45in) {
{%- endraw %}
            {{ cookiecutter.short_course_name }}
{%- raw %}
        };
    \node[bubble sheet fill in,anchor=south west]
        at (0.5in,-1in) {
{%- endraw %}
            {{ cookiecutter.get('section_number', '') }}
{%- raw %}
        };
    \node[bubble sheet fill in,
        anchor=south west]
        at (0.5in,-1.5in) {
            \makeatletter\@date\makeatother
        };
\end{tikzpicture}
\examcleardoublepage
\fi
{%- endraw -%}
{%- endif -%}
{%- raw %}
\end{coverpages}

\begin{EnvFullwidth}
\section*{Multiple Choice}
{%- endraw -%}
{% if cookiecutter.use_bubble_sheet -%}
{%- raw %}
\begin{instructions}%
    Shade the circles corresponding to your answers to the multiple choice
    questions on the bubble sheet. 
%%    Pages \pageref{begin-fixed-response}--\pageref{end-fixed-response} will not be read.
    No explanation is required.
\end{instructions}
{%- endraw -%}
{% else -%}
{%- raw %}
\begin{instructions}%
    Shade the circles corresponding to your answers to the multiple choice
    questions, as in the example below:
    \begin{center}
    \begin{minipage}{0.95\textwidth} 
        \begin{itemize}
            \item[\samplecirc{A}] Some answer
            \item[\samplecircSoln{B}] The answer you think is correct
            \item[\samplecirc{C}] Some answer
            \item[\samplecirc{D}] Some answer
            \item[\samplecirc{E}] Some answer
        \end{itemize}
    \end{minipage}  
    \end{center}
    \medskip
    No explanation is required to be shown and no partial credit will be given
    for the multiple choice questions.
\end{instructions}
{%- endraw -%}
{% endif -%}
{%- raw %}
\end{EnvFullwidth}


\begin{questions}
    % \vspace*{-0.6in}

{% endraw -%}
{%- set num_mcq = cookiecutter.num_mcq | int -%}
{%- if (num_mcq > 0) -%}    
    \begingradingrange{mc}
{% for i in range(0,num_mcq) %}    
    \begin{question}[{{ cookiecutter.points_per_mcq }}]
        Stem
        \begin{choices}
            \CorrectChoice{}
            \choice{}
            \choice{}
            \choice{}
            \choice{}
        \end{choices}
    \end{question}
{% endfor %}
    \endgradingrange{mc}
    \examcleardoublepage%
{% endif %}

{%- set num_cata = cookiecutter.num_cata | int -%}
{%- if (num_cata > 0) -%}
\begin{EnvFullwidth}
    \section*{Check All That Apply Groups}
    \begin{instructions}%
        In each group below, more than one option \textbf{may} be correct.
        Each option is its own question on the bubble sheet:
        fill in \textbf{A} (Yes) if the option should be chosen, and \textbf{B} (No) otherwise.
        Each option is worth 1 point, so each group is worth 5 points.
    \end{instructions}
\end{EnvFullwidth}

    \begingroup
    \begingradingrange{cata}
    \makeatletter\let\point@block\relax\makeatother
    \NewList{cataoptions}
    \newcommand{\Yes}{\hfill\mbox{\begin{oneparchoices}[norandomize]\CorrectChoice{Yes}\choice{No}\end{oneparchoices}}}
    \newcommand{\No}{\hfill\mbox{\begin{oneparchoices}[norandomize]\choice{Yes}\CorrectChoice{No}\end{oneparchoices}}}
{% for i in range(1,num_cata+1) %}
    \begingradingrange{{ ("msq-" ~ i) | embrace }}
    \begin{EnvUplevel}
        \firstqinrange{{ ("msq-" ~ i) | embrace }}--\lastqinrange{{ ("msq-" ~ i) | embrace }}.
        Stem: which of the following \dots?
    \end{EnvUplevel}
    \SetList{cataoptions}{
      {Option}{\Yes},
      {Option}{\Yes},
      {Option}{\No},
      {Option}{\No},
      {Option}{\No}
    }
    \ForEachRandomItem{cataoptions}{Option}{
        \ReadFieldItem{\Option}{0}{Stem}%
        \ReadFieldItem{\Option}{1}{Answer}
        \begin{question}[\mspoints]\Stem\Answer\end{question}
    }
    \endgradingrange{{ ("msq-" ~ i) | embrace }}
    \medskip
{% endfor %}
    \endgradingrange{cata}
    \endgroup
    \examcleardoublepage%
{% endif %}

{%- set num_tf = cookiecutter.num_tf | int -%}
{%- if (num_tf > 0) -%}
    \begingradingrange{tfj}
    \begingroup
    \makeatletter\let\point@block\relax\makeatother

\section*{True/False with Reasons}
\begin{EnvFullwidth}
\begin{instructions}
    Decide if the statement is true or false, and give a reason.
    Each statement on the exam counts for two questions on the bubble sheet.
    \begin{itemize}
        \item For the first, select \textbf{A} if the statement is true and \textbf{B} if it is false.
          This question is worth 1 point.
        \item For the second, select a reason.
        \begin{itemize}
            \item Select \textbf{A} if the statement is true by \textbf{definition}
            \item Select \textbf{B} if the statement has been or can be \textbf{proven} 
            \item Select \textbf{C} if the statement is \textbf{vacuously} true
            \item Select \textbf{D} if the statement is false, because it is only \textbf{sometimes} satisfied
            \item Select \textbf{E} if the statement is false, because it is \textbf{never} satisfied
        \end{itemize}
        This question is worth 2 points.
    \end{itemize}
\end{instructions}
\end{EnvFullwidth}

    \NewList{tfjqs}
{% for i in range(1,num_tf+1) %}
    \InsertLastItem{tfjqs}{{ bgroup }}%
        \begingradingrange{{ ("tfj-" ~ i) | embrace }}
        \begin{EnvUplevel}
            \firstqinrange{{ ("tfj-" ~ i) | embrace }}--\lastqinrange{{ ("tfj-" ~ i) | embrace }}.
            Statement
        \end{EnvUplevel}
        \begin{question}[\tfpoints] Truth value
            \begin{oneparchoices}[norandomize]
                \CorrectChoice{True}
                \choice{False}
            \end{oneparchoices}
        \end{question}
        \begin{question}[\tfrpoints] Reason
            \begin{oneparchoices}[norandomize]
                \CorrectChoice{Definition}
                \choice{Proven}
                \choice{Vacuous}
                \choice{Sometimes}
                \choice{Never}
            \end{oneparchoices}
            \begin{solution}
                Reason
            \end{solution}
        \end{question}
        \endgradingrange{{ ("tfj-" ~ i) | embrace }}
    }
{% endfor %}
    \ForEachRandomItem{tfjqs}{Question}{
        \Question
        \medskip
    }
    \endgroup
    \endgradingrange{tfj}
    \examcleardoublepage%
{% endif %}
{%- raw %}
\begingradingrange{fr}%
    \fullwidth{\section*{Free Response}

    \begin{instructions}%
        The remaining questions are free response questions or proofs.  Put your
        answers in the blanks (where provided) and your work/explanations in the
        space below each problem or problem part.
        Justify your steps with statements proven in the textbook or notes.
        Write complete sentences in English. Do not use the three-dots notation (\(\therefore\) or \(\because\)).
        Command terms such as \textbf{explain} and \textbf{prove} are defined in the glossary
        on page~\pageref{command-verbs-glossary}.
    \end{instructions}}
{% endraw %}
{%- set num_frq = cookiecutter.num_frq | int -%}
{% for i in range(0,num_frq) %}
    \begin{question}[10]
        Stem. \Gls{prove} \dots
        \begin{solution}[\stretch{1}]
            Solution
        \end{solution}
    \end{question}
    \examcleardoublepage
{% endfor %}
\endgradingrange{fr}
{%- raw %}

    \end{questions}
{%- endraw %}
{%- if cookiecutter.get('use_formula_sheet') %}
{%- raw %}
% The formula sheet is the front and back of the last sheet of paper:
% \examcleardoublepage puts us at the top of an odd page.
\ifprintanswers\else
    \examcleardoublepage
    \includepdf[pages=-,noautoscale,scale=1.15]{formula_sheet_nyu}
\fi
{%- endraw %}
{%- endif %}
{%- raw %}
\end{document}
%</questions>
{% endraw -%}